% Generated draft: requires mathematical and source review.
\documentclass[11pt]{article}
\usepackage[margin=25mm]{geometry}
\usepackage{fontspec}
\setmainfont{Latin Modern Roman}
\usepackage{amsmath,amssymb,mathtools,unicode-math}
\setmathfont{Latin Modern Math}
\usepackage{enumitem,needspace,xurl,hyperref}
\setlength{\parindent}{0pt}
\setlength{\parskip}{5pt}
\setlength{\emergencystretch}{3em}
\newcommand{\N}{\mathbb{N}}
\newcommand{\Z}{\mathbb{Z}}
\newcommand{\Q}{\mathbb{Q}}
\newcommand{\R}{\mathbb{R}}
\newcommand{\C}{\mathbb{C}}
\newcommand{\F}{\mathbb{F}}
\newcommand{\A}{\mathbb{A}}
\newcommand{\PP}{\mathbb{P}}
\newcommand{\E}{\mathbb{E}}
\newcommand{\eps}{\varepsilon}
\newcommand{\dd}{\,\mathrm{d}}
\begin{document}
\section*{UnsolvedMath ID: 10400173}
\textbf{Source problem code:} catalogue entry
\subsection{Definitions and mathematical statement}
A subfactor is an inclusion \(P\subset Q\) of von Neumann factors. The source asks for one whose associated three-manifold theory distinguishes \(L(7,1)\) from \(L(7,2)\). In a fixed lens-space convention,
\[
 L(p,q)=S^3/\bigl((z_1,z_2)\sim
 (e^{2\pi i/p}z_1,e^{2\pi iq/p}z_2)\bigr),\qquad\gcd(p,q)=1.
\]
If \(Z_{P\subset Q}\) denotes the intended associated invariant, the concrete target is
\[
 Z_{P\subset Q}(L(7,1))\ne Z_{P\subset Q}(L(7,2)).
\]
The subfactor type/index/depth, invariant construction, normalization and orientation conventions are not supplied, so this notation remains a placeholder for the source's intended theory.

The broader request is to find subfactor theories classifying three-manifolds ``as well as possible''. It does not prescribe a completeness criterion, comparison class or optimization ordering; no complete all-three-manifold classifier is silently asserted as the required conclusion.

The passage cites the Ponzano--Regge \(\operatorname{SU}(2)\) state-sum model and its asymptotics as motivation. It suggests that some manageable infinite-depth, strongly amenable subfactors might have enough growth control to handle analogous asymptotic state sums. This is expressly tentative. Strong amenability is a technical condition not defined by the excerpt's informal growth description, and convergence, regularization, triangulation independence and asymptotic control of an infinite sum would need exact hypotheses. None is inferred simply from the word amenable.

The supplied summary gives only an open/triage attribution and no distinguishing subfactor or current theorem. The historical concrete target, broader classification request and tentative asymptotic proposal are retained separately.
\subsection{Short English statement}
Find a subfactor-based three-manifold invariant separating \(L(7,1)\) and \(L(7,2)\), then seek broader classification power. The suggested use of strongly amenable infinite-depth subfactors for state-sum asymptotics remains a separate tentative direction.
\subsection{Sources}
\begin{itemize}
\item Local registry: \url{data/unsolvedmath/problems.json}, permanent UnsolvedMath ID \(10400173\), source code catalogue entry; supplied statement, background, and research summary, updated 2026-08-06.
\item Ohtsuki, \emph{Problems on invariants of knots and 3-manifolds} (2002), Problem 9.9, attributed to N. Sato, PDF page 141: \url{https://msp.org/gtm/2002/04/gtm-2002-04-024s.pdf}. Supplied PDF-font-structure extraction accessed 2026-07-29.
\item Ponzano--Regge, source reference [332], cited for the \(\operatorname{SU}(2)\) state-sum model and asymptotics; its full bibliographic entry is not included.
\end{itemize}
\end{document}
